%! Designing Experiments — Unit 3, Lesson 3.4 — Warm-Up (Answer Key)
% Mirrors the blank exactly apart from the package swap, the pageheader, and the answers.
% NO teachernote here — teacher prose lives in the lesson plan.
% ONE PAGE, blank and key. Three items; do not add a fourth.
% GRADUAL RELEASE: the warm-up activates prior knowledge before the guided notes.
%   Item 1 — spiral to Lesson 3.1: observational study or experiment, and why an observed
%            association (breakfast and GPA) is not a cause. Seeds CONFOUNDING VARIABLE.
%   Item 2 — a review game used in 1st period and not 6th: name another reason 1st period
%            scored higher. Seeds confounding, COMPARISON, and why chance should assign.
%   Item 3 — split 20 students into two groups of 10 by chance with slips of paper. Seeds
%            RANDOM ASSIGNMENT (EK VAR-3.C.2) — the move every design in the notes uses.
% Every display is given; nothing is sketched.
\documentclass[12pt]{article}
\usepackage{apstats-article}
\usepackage{apstats-key}   % pulls in -boxes; do NOT also load -boxes

\begin{document}

\pageheader{Unit 3, Lesson 3.4}{Warm-Up --- Answer Key}
% NAMESTRIP: no \namedateperiod here — the name row lives on the cover only.

\vspace{0.28in}

\begin{enumerate}[leftmargin=1.9em, itemsep=0.42in]

  \item \textbf{From Lesson 3.1.} A survey of 500 teens finds that those who eat breakfast
        every day have higher GPAs than those who do not.

        \vspace{12pt}
        Observational study or experiment? \ans{Observational study}
        \vspace{16pt}

        Does this show that eating breakfast \emph{causes} a higher GPA? Why or why not?
        \vspace{12pt}

        \ansline{No. Nothing was imposed. Another variable, such as a steady home routine,}\\[18pt]\ansline{could lead to both eating breakfast and a higher GPA.}

  \item A teacher wants to know whether a review game raises quiz scores. She plays the game
        with her 1st-period class but not her 6th-period class, and 1st period scores higher on
        the quiz. Give one \emph{other} reason 1st period might have scored higher.

        \vspace{12pt}
        \ansline{Time of day: students may be more alert in the morning. Or 1st period may}\\[18pt]\ansline{simply have stronger students (an honors section, say).}

  \item The teacher tries again with 20 students in one class. Describe how she could use 20
        slips of paper to split the students into two groups of 10 \emph{by chance}, one group
        to play the game and one not.

        \vspace{12pt}
        \ansline{Write each student's name on a slip, mix the slips in a hat, and draw 10}\\[18pt]\ansline{without replacement. Those 10 play the game; the other 10 do not.}\\[18pt]\ansline{}

\end{enumerate}

\end{document}
